% !PRINTED 484
\weberpage{484}
%<<CONT
hence:
\[ f_2\left(\frac{\omega}{n}\right) = e^{-\frac{r\pi i}{24}}x, \]
\[ f_2(n\omega) = \frac{e^{-\frac{r\pi i}{24}}\sqrt2}{x}, \]
if, according as $r$ is even or odd,
\weq{33}{x = f_1(\sqrt{-m}) \quad\text{or}\quad = f(\sqrt{-m})}
% [sic] German source omits $x$ before the second equality
is set. If, then, in the modular equation belonging to $n$,
 II., § 73, we take:
\[ u_1 = f_2(\omega), \]
\[ v_1 = \left(\frac2n\right) f_2(n\omega) = \left(\frac2n\right) \frac{e^{-\frac{r\pi i}{24}}\sqrt2}{x} \]
or
\[ v_1 = f_2\left(\frac{\omega}{n}\right) = e^{-\frac{r\pi i}{24}}x, \]
there result two equations, from which by elimination
 of $u_1$ one derives an equation for $x$.
In order to carry out this calculation for $n = 5$, we set:
\weq{34}{\begin{aligned}
x f_2(\omega) &= \sqrt2\,e^{-\frac{5r\pi i}{24}}\xi, \\
\frac{x}{f_2(\omega)} &= e^{-\frac{5r\pi i}{24}}\eta,
\end{aligned}}
from which, corresponding to the two assumptions for $v_1$, one obtains
for an odd $r$ the equations:
\weq{35}{\begin{aligned}
\xi^3+\frac1{\xi^3}+2\left(\eta^2-\frac1{\eta^2}\right)&=0, \\
\eta^3+\frac1{\eta^3}+2\left(\xi^2-\frac1{\xi^2}\right)&=0.
\end{aligned}}
These equations hold for $r = 1, 3, 7$, which belong to the values
$m = 49, 41, 1$.
If one subtracts the two equations (35), one can cast off
the factor $\xi - \eta$, which can vanish only for $m = 1$:
\[ [(\xi+\eta)^2 - \xi\eta]\left(1-\frac1{\xi^3\eta^3}\right) - 2(\xi+\eta)\left(1+\frac1{\xi^2\eta^2}\right) = 0, \]
%CONT>>
